\documentclass[11pt]{article}
\usepackage[T1]{fontenc}
\usepackage{lmodern,amsmath,amssymb,amsthm,mathtools}
\usepackage[margin=1in]{geometry}
\usepackage{microtype}
\usepackage[colorlinks=true,linkcolor=blue,citecolor=blue,urlcolor=blue]{hyperref}
\newtheorem{theorem}{Theorem}
\newtheorem{proposition}[theorem]{Proposition}
\newcommand{\Z}{\mathbb Z}
\newcommand{\R}{\mathbb R}
\newcommand{\dd}{\,\mathrm d}
\DeclareMathOperator{\Var}{V}
\title{Bilinear ergodic Hilbert variation for commuting transformations\\as a consequence of annular triangular variation}
\author{Alec Kriebel\\\small\href{https://orcid.org/0009-0001-9320-500X}{ORCID: 0009-0001-9320-500X}}
\date{October 6, 2026}
\hypersetup{pdftitle={Bilinear ergodic Hilbert variation for commuting transformations as a consequence of annular triangular variation},pdfauthor={Alec Kriebel}}
\begin{document}
\maketitle
\begin{abstract}
We give an explicit restriction and transference consequence of OpenAI's
October 5, 2026 full annular variation theorem for the triangular Hilbert
transform. For arbitrary commuting invertible probability-preserving
transformations $S,T$ and complex $f,g\in L^3$, the symmetric bilinear
Hilbert sums have bounded pointwise $r$-variation in $L^{3/2}$ for every
$r>2$, and converge almost everywhere and in $L^{3/2}$. The restriction
uses boxes of fixed positive area and a uniformly summable kernel error.
The analytic breakthrough is the cited continuous theorem; the present
note supplies the discrete restriction and full measure-theoretic transfer.
\end{abstract}

\section{Statement and provenance}
Let $(X,\Sigma,\mu)$ be a probability space. Assume that $S,T$ are
invertible measurable measure-preserving transformations with measurable
inverses, and $ST=TS$. For complex $f,g\in L^3(X)$ put
\[
 H_0(f,g)=0,\qquad
 H_N(f,g)(x)=\sum_{0<|n|\le N}\frac{f(S^n x)g(T^n x)}n\quad(N\ge1).
\]
For a complex sequence $(a_N)_{N\ge0}$ define the homogeneous variation
\begin{equation}\label{eq:variation}
 \Var^r(a_N:N\ge0)
 =\sup_{\substack{J\ge1\\0\le N_0<\cdots<N_J\,,\ N_j\in\Z}}
 \left(\sum_{j=1}^J|a_{N_j}-a_{N_{j-1}}|^r\right)^{1/r}.
\end{equation}
If $a_0=0$, prepending $0$ to any list shows that restricting the
supremum to $N_0=0$ gives the same value. In particular
$\sup_N|a_N|\le\Var^r(a)$.

\begin{theorem}\label{thm:ergodic}
For every $r>2$ there is a constant $D_r<\infty$, independent of the
probability space and transformations, such that
\begin{equation}\label{eq:ergodic-var}
 \|\Var^r(H_N(f,g):N\ge0)\|_{3/2}
 \le D_r\|f\|_3\|g\|_3.
\end{equation}
The maximal function satisfies the same bound with $\Var^r$ replaced
by $\sup_{N\ge0}|H_N|$. The sums converge almost everywhere and in
$L^{3/2}$ to a complex bilinear operator $H(f,g)$, and moreover
\begin{equation}\label{eq:tail}
 \left\|\sup_{N\ge m}|H_N(f,g)-H(f,g)|\right\|_{3/2}\longrightarrow0.
\end{equation}
\end{theorem}

The single-transformation Hilbert theorem of Demeter
\cite[Theorem 1.2]{Demeter} treats opposite powers of one
transformation. Becker and Durcik \cite[Corollary 1.4]{BD} prove sharp
$r>2$ variation for bilinear averages along opposite powers of one
transformation; their discussion of equation (1.8) distinguishes the
arbitrary commuting case. Demeter and Thiele \cite[Section 6]{DT}
relate the triangular transform to two commuting transformations.
Calder\'on's transference method \cite{Calderon} is established machinery;
we include its finite-window argument rather than appeal to it without
quantifiers.
The continuous-to-lattice and orbit-array reductions already occur in
\cite[Section 6, footnote 36]{DT}; positive-area phase integration also
appears in \cite[Section 5]{DKST}. The fixed-width singular-kernel
bookkeeping below makes this particular consequence explicit.

The new input is the following continuous estimate, stated as
Theorem 1.1 in the October 5 annular manuscript \cite{OAI}. For
$F,G\in L^3(\R^2)$ write
\[
 B_{a,b}(F,G)(u,v)=\int_{a<|t|<b}F(u+t,v)G(u,v+t)\,\frac{\dd t}{t}.
\]
Then, for every $r>2$,
\begin{equation}\label{eq:continuous}
 \left\|\sup_{\substack{J\ge1\\0<t_0<\cdots<t_J\,,\ t_j\in\mathbb Q}}
       \left(\sum_{j=1}^J|B_{t_{j-1},t_j}(F,G)|^r\right)^{1/r}
       \right\|_{L^{3/2}(\R^2)}
 \le C_r\|F\|_3\|G\|_3.
\end{equation}
The supremum is inside the norm and partitions may depend on $(u,v)$.
It includes all positive scales, with unrestricted finite partition
length. We use precisely this estimate, rather than a maximal, dyadic,
or fixed-partition norm-variation estimate. The upstream theorem is
not reproved here; all conclusions below are consequences of it.
The listed family 082 Lean scope concerns an older maximal theorem
and does not formally verify \eqref{eq:continuous} or this note.

Our contribution is the fixed-width restriction proof in Section 2
and its explicit transfer in Section 3. We do not claim independent
resolution of the continuous triangular Hilbert problem or invention
of transference. Manuscript dates in the upstream collection are not
assertions of first public disclosure. The source repository version used is
\[
\texttt{adc7f1241b42e322a6451854ab7e4b4c146bf78a}.
\]

\section{Fixed-width restriction to the lattice}
For finitely supported arrays $A,B:\Z^2\to\mathbb C$ put
\[
 h_N(A,B)(i,j)=\sum_{0<|k|\le N}\frac{A(i+k,j)B(i,j+k)}k,
 \qquad h_0=0.
\]
All lattice norms use counting measure.

\begin{proposition}\label{prop:discrete}
Assuming \eqref{eq:continuous},
\begin{equation}\label{eq:discrete}
 \|\Var^r(h_N:N\ge0)\|_{\ell^{3/2}(\Z^2)}
 \le D_r\|A\|_{\ell^3}\|B\|_{\ell^3},
 \qquad D_r=8\,4^{2/3}C_r+\frac{\pi^2}{13}.
\end{equation}
\end{proposition}

\begin{proof}
Fix $h=1/8$ and $q=1/16$. Extend $A,B$ to $F,G$ on $\R^2$ by
making them constant with their corresponding array values on
\[
 (i,j)+[-h,h]^2\quad((i,j)\in\Z^2)
\]
and zero elsewhere. Boundary values have no effect and can be assigned
arbitrarily. The boxes are disjoint and
\begin{equation}\label{eq:input-norm}
 \|F\|_3=(4h^2)^{1/3}\|A\|_{\ell^3},\qquad
 \|G\|_3=(4h^2)^{1/3}\|B\|_{\ell^3}.
\end{equation}
Evaluate at $(u,v)=(i+a,j+b)$ with $a,b\in(-q,q)$.
If $F(u+t,v)G(u,v+t)$ is nonzero, the fixed coordinate $v$ selects
the second cell index $j$ of $F$, and the fixed coordinate $u$ selects
the first cell index $i$ of $G$. The moving coordinates require
\[
 t\in k+[-h-a,h-a]\quad\hbox{and}\quad
 t\in l+[-h-b,h-b]
\]
for integers $k,l$. If these intervals overlap, then
$|k-l|\le2h+|a-b|\le3/8<1$, so $k=l$.
Consequently the integrand equals the single lattice product
$A(i+k,j)B(i,j+k)$ on $k+I_{a,b}$, where
\begin{equation}\label{eq:overlap}
 I_{a,b}=[-h-\min(a,b),\ h-\max(a,b)],\qquad
 \ell(a,b)=|I_{a,b}|=2h-|a-b|\ge1/8.
\end{equation}
Also $\ell\le1/4$ and every $s\in I_{a,b}$ satisfies $|s|\le3/16$.
Different integer windows are disjoint.

Set $\rho_N=N+1/2$ for $N\ge0$. For $k\ne0$, the entire window
$k+I_{a,b}$ lies in $|t|\in[|k|-3/16,|k|+3/16]$.
Thus $\rho_M<|t|<\rho_N$ includes that window exactly when
$M<|k|\le N$. This verifies negative as well as positive windows.
The window at $k=0$ lies in $|t|\le3/16$ and is excluded even at
the smallest endpoint $\rho_0=1/2$. Hence for $0\le M<N$,
\begin{equation}\label{eq:kernel-identity}
 \begin{split}
 B_{\rho_M,\rho_N}(F,G)(i+a,j+b)
 &=\sum_{M<|k|\le N}A(i+k,j)B(i,j+k)K_k(a,b),\\
 K_k(a,b)&=\int_{I_{a,b}}\frac{\dd s}{k+s}.
 \end{split}
\end{equation}
There is no principal value in this identity. In particular $k+s$
has the sign of $k$.

Writing $K_k=\ell/k+\delta_k$, the exact identity
$1/(k+s)-1/k=-s/[k(k+s)]$ gives the uniform bound
\begin{equation}\label{eq:error}
 \frac{|\delta_k(a,b)|}{\ell(a,b)}
 \le\frac{3}{13}\frac1{k^2}\qquad(k\ne0),
\end{equation}
since $|k+s|\ge(13/16)|k|$.
Let $E_N=\sum_{0<|k|\le N}\delta_k A(i+k,j)B(i,j+k)$ and $E_0=0$.
For any increasing integer list, the radial shells of its increments
are disjoint. The $\ell^r$ norm of the error increments is at most their
$\ell^1$ norm, which is at most the sum of all absolute error terms.
Therefore
\[
 \ell^{-1}\Var^r(E_N:N\ge0)
 \le e(i,j):=\frac3{13}\sum_{k\ne0}
                  \frac{|A(i+k,j)B(i,j+k)|}{k^2}.
\]
Minkowski, H\"older, and translation invariance imply
\begin{equation}\label{eq:error-norm}
 \|e\|_{\ell^{3/2}}
 \le\frac3{13}\sum_{k\ne0}\frac1{k^2}
                      \|A\|_{\ell^3}\|B\|_{\ell^3}
 =\frac{\pi^2}{13}\|A\|_{\ell^3}\|B\|_{\ell^3}.
\end{equation}
This bounds the whole error variation, not merely each truncation.

Let $W(F,G)$ denote the pointwise supremum in
\eqref{eq:continuous}. The half-integer endpoints are rational.
Taking the supremum over integer lists in \eqref{eq:kernel-identity},
using the variation triangle inequality and $\ell^{-1}\le8$, gives
\begin{equation}\label{eq:boxwise}
 \Var^r(h_N)(i,j)\le8W(F,G)(i+a,j+b)+e(i,j)
\end{equation}
for almost every $(a,b)\in(-q,q)^2$, for every $(i,j)$.
All lists are countable, so one common exceptional set suffices.
Apply the $L^{3/2}$ triangle inequality on the disjoint output boxes,
whose area is $4q^2=1/64$. The left side is constant on each box,
so \eqref{eq:boxwise}, \eqref{eq:continuous} and
\eqref{eq:input-norm} give
\[
 \|\Var^r(h_N)\|_{\ell^{3/2}}
 \le8(4q^2)^{-2/3}C_r(4h^2)^{2/3}
                        \|A\|_{\ell^3}\|B\|_{\ell^3}
      +\|e\|_{\ell^{3/2}}.
\]
The coefficient equals $8\,4^{2/3}C_r$, proving the proposition.
The restriction is carried by a fixed set of positive area;
no evaluation on a null set of suspension phases is used.
\end{proof}

\section{Finite-window transference and convergence}
\begin{proof}[Proof of Theorem \ref{thm:ergodic}]
Choose measurable representatives of $f,g$ taking finite values outside
their null sets. Remove all preimages of those null sets under finite
words in $S,S^{-1},T,T^{-1}$.
The resulting invariant full-measure set carries every value in the
finite sums. A change of representatives affects only such a countable
union of null sets. The same convention treats transformations commuting
modulo null sets: remove all such word preimages of the failure set, after
which all required identities hold on the invariant complement.

Fix an integer $M\ge1$. Let $W_M(x)$ be the maximum in
\eqref{eq:variation} over lists with endpoints in $\{0,\ldots,M\}$.
There are finitely many such lists, so $W_M$ is measurable. For
$L\ge1$ let $Q_L=\{-L,\ldots,L\}^2$. At each $x$ form finite arrays
\[
 \begin{split}
 A_x(i,j)&=f(S^iT^j x)\mathbf1_{Q_{L+M}}(i,j),\\
 B_x(i,j)&=g(S^iT^j x)\mathbf1_{Q_{L+M}}(i,j).
 \end{split}
\]
For $(i,j)\in Q_L$ and $|n|\le M$, every required entry is in
$Q_{L+M}$. Commutation gives exactly
\[
 h_N(A_x,B_x)(i,j)=H_N(f,g)(S^iT^j x)\quad(0\le N\le M).
\]
Proposition \ref{prop:discrete}, with $p=3/2$, implies
\[
 \sum_{(i,j)\in Q_L}W_M(S^iT^j x)^p
 \le D_r^p
 \left(\sum_{Q_{L+M}}|f(S^iT^j x)|^3\right)^{1/2}
 \left(\sum_{Q_{L+M}}|g(S^iT^j x)|^3\right)^{1/2}.
\]
Integrating, measure invariance on the left and Cauchy--Schwarz on
the right give
\[
 |Q_L|\|W_M\|_p^p
 \le D_r^p|Q_{L+M}|\|f\|_3^{3/2}\|g\|_3^{3/2}.
\]
The finite right side establishes integrability without assuming it in
advance. For fixed $M$, let $L\to\infty$;
$|Q_{L+M}|/|Q_L|\to1$. Finally $W_M$ increases to the full countable
supremum \eqref{eq:variation}. Monotone convergence proves
\eqref{eq:ergodic-var}. Partitions may depend on the output throughout.

Fix, for example, $r=3$. The variation is finite almost everywhere.
Any scalar sequence with finite $r$-variation is Cauchy: otherwise
some $\varepsilon>0$ admits arbitrarily late pairs of indices whose
values differ by at least $\varepsilon$. Recursively choose $J$
strictly ordered disjoint pairs. A partition containing their endpoints
has variation at least $\varepsilon J^{1/r}$, a contradiction.
Completeness of $\mathbb C$ gives a finite limit $H$. It is measurable
as the limit along integers, setting it to zero on the exceptional set.

Since $H_0=0$, $H_*:=\sup_N|H_N|\le\Var^3(H_N)$ belongs to $L^p$.
Also $|H|\le H_*$. The measurable tail
$\sup_{N\ge m}|H_N-H|$ tends to zero pointwise and is bounded by
$2H_*$. Dominated convergence applied to its $p$th power proves
\eqref{eq:tail}, hence norm convergence. The maximal bound follows
from \eqref{eq:ergodic-var}; the limit bound follows from domination.
Norm limits of the bilinear identities for $H_N$ show that $H$ is
complex bilinear on $L^3\times L^3$.
\end{proof}

\section{Scope and verification disclosure}
The truncations are symmetric and the kernel is odd. This proof gives
no one-sided commuting Ces\`aro convergence, no noncommuting theorem,
no $r=2$ estimate, and no extension of the input exponents beyond $L^3$.
No numerical evidence substitutes for the analytic estimate.
Supplemental exact kernel checks test only Section 2 bookkeeping.

The continuous theorem is attributed to OpenAI using its supplied
manuscript citation. The dependency ledger records the version and
scope of independent automated proof audits. Such audits are evidence,
not formal certification. No theorem in this note is claimed to have
been formalized. AI tools were used extensively in research, drafting
and verification. This preprint has not undergone conventional human
peer review or refereeing at publication.

\begin{thebibliography}{9}
\bibitem{BD} L. Becker and P. Durcik,
\emph{The shifted bilinear Hilbert transform}, arXiv:2603.20173v1 (2026),
\url{https://arxiv.org/abs/2603.20173}.
\bibitem{Calderon} A. P. Calder\'on,
\emph{Ergodic theory and translation-invariant operators},
Proc. Natl. Acad. Sci. U.S.A. \textbf{59} (1968), 349--353,
\href{https://doi.org/10.1073/pnas.59.2.349}{doi:10.1073/pnas.59.2.349}.
\bibitem{Demeter} C. Demeter,
\emph{Pointwise convergence of the ergodic bilinear Hilbert transform},
arXiv:math/0601277v1 (2006),
\url{https://arxiv.org/abs/math/0601277}.
\bibitem{DT} C. Demeter and C. Thiele,
\emph{On the two dimensional bilinear Hilbert transform},
Amer. J. Math. \textbf{132} (2010), 201--256,
\url{https://arxiv.org/abs/0803.1268}.
\bibitem{DKST} P. Durcik, V. Kova\v{c}, K. A. \v{S}kreb and C. Thiele,
\emph{Norm-variation of ergodic averages with respect to two commuting
transformations}, Ergodic Theory Dynam. Systems \textbf{39} (2019),
658--688, \url{https://arxiv.org/abs/1603.00631}.
\bibitem{OAI} OpenAI,
\emph{Annular variation of the triangular Hilbert transform at the
symmetric point}, OpenAI Math Release preprint (manuscript dated
October 5, 2026), source version
\texttt{adc7f1241b42e322a6451854ab7e4b4c146bf78a},
\href{https://github.com/openai/math/blob/adc7f1241b42e322a6451854ab7e4b4c146bf78a/preprints/Annular-variation-of-the-triangular-Hilbert-transform-at-the-symmetric-point-October-5-2026/annular-variation.pdf}{manuscript and source repository}.
\end{thebibliography}
\end{document}
